\begin{solapartat}
La força de Lorentz ve donada per l'expressió:
\[ \vec{F} = q\,\vec{v}\wedge\vec{B} \quad\punts{0,2} = 2\cdot 1,60\cdot 10^{-19}\cdot 8\cdot 10^{5}\cdot 1,2\cdot(\vec{\imath}\wedge\vec{k}) = -3,07\cdot 10^{-13}\,\vec{\jmath}\ \mathrm{N} \quad\punts{0,2} \]
La trajectòria serà circular ja que la força és perpendicular a la la velocitat de la partícula, % DUBTE: errata «a la la» a la pauta original
 \punts{0,2} i girarà en el pla $xy$ en sentit horari. \punts{0,1}

\figura[0.3]{sol_fig1.png}
\aladreta{\punts{0,3}}
\end{solapartat}

\begin{solapartat}
La força de Lorentz serà la força centrípeta que ens farà girar la partícula $\alpha$: \punts{0,2}
\[ m_\alpha\,\frac{v^2}{r} = q\,v\,B \quad\punts{0,2} \Rightarrow r = \frac{m_\alpha\,v}{q\,B} = 1,38\cdot 10^{-2}\ \mathrm{m} \quad\punts{0,2} \Rightarrow \]
\[ 2\pi\,\nu\,r = v \quad\punts{0,2} \Rightarrow \nu = \frac{v}{2\pi\,r} = \frac{qB}{2\pi\,m_\alpha} = 9,20\cdot 10^{6}\ \mathrm{Hz} = 9,2\ \mathrm{MHz} \quad\punts{0,2} \]
\end{solapartat}
